(1) 时间 \(dt\) 内入料质量 \(\alpha dt\)、携带水平动量 \(u\alpha dt\)，车原质量 \(m(t)=m_0+\alpha t\)。将小车与即将捕获的落料作为短时整体，动量守恒及外力冲量给 \(d(mv)/dt=F+\alpha u\)，即 \(m\dot v=F+\alpha(u-v)\)；漏掉 \(-\alpha v\) 会错误处理随车增加的惯量。
(2) 积分得到 \(m(t)v(t)=m_0v_0+(F+\alpha u)t\)，从而令 \(v_* =u+F/\alpha\)，有 \(v(t)=v_*+(v_0-v_*)m_0/(m_0+\alpha t)\)。再积分得 \(x(t)-x(0)=v_*t+(v_0-v_*)m_0\ln(1+\alpha t/m_0)/\alpha\)，将 \(t=0\) 代入可同时复现初始速度和位置。
(3) 小车动能 \(K=mv^2/2\)，由 \(m\dot v=F+\alpha(u-v)\) 计算 \(\dot K=Fv+\alpha uv-\alpha v^2/2\)。入料动能流为 \(\alpha u^2/2\)，两者相减给吸附耗散率 \(\dot Q_{\rm loss}=\alpha(u-v)^2/2\ge0\)，因此 \(\dot K=Fv+\alpha u^2/2-\dot Q_{\rm loss}\)，功率单位与能量收支一致。
(4) 无外力时 \(mv=m_0v_0+\alpha ut\)，小车最终趋近来料速度 \(u\)；若 \(u=0\) 则水平动量 \(mv=m_0v_0+Ft\)。在 \(t\) 固定下 \(\alpha\to0\)，对数展开及速度分式均给 \(v(t)\to v_0+Ft/m_0\)、\(x-x_0\to v_0t+Ft^2/(2m_0)\)。只有 \(u=v\) 时瞬间相对碰撞损耗为零，否则完全粘附属于非弹性过程。
